
\subsubsection{Structured Error Format}

The structured format transforms raw compiler output into a consistent template:

\begin{verbatim}
=== ERROR REPORT ===

PROBLEM: [error_type]: [error_message]

LOCATION: Line [line_number] in [file_name]

CONTEXT:
[surrounding_code_snippet]
    ^ [pointer_to_error]

ROOT_CAUSE: [explanation_of_why_error_occurred]

FIX_HINT: [optional, level-dependent guidance]
\end{verbatim}

The section labels create parsing anchors for identifying relevant information. The consistent structure enables pattern matching against similar examples in training data.

\subsubsection{Error Parser Implementation}

A parser extracts structured information from Python tracebacks using regex-based extraction. The implementation supports eight common error types: NameError, TypeError, IndexError, KeyError, ZeroDivisionError, ValueError, AttributeError, and SyntaxError. The parser outputs a StructuredError dataclass containing line\_number, error\_type, error\_message, and code\_context fields.

\subsubsection{Experimental Conditions}

\begin{table}[htbp]
\centering
\resizebox{\columnwidth}{!}{%
\begin{tabular}{lll}
\toprule
Condition & Description & Purpose \\
\midrule
Raw & Direct compiler/linter output & Baseline \\
Structured & Template format with section labels & Treatment \\
Scrambled & Structured content, random section order & Information control \\
\bottomrule
\end{tabular}
}
\end{table}

The Scrambled condition contains identical information to Structured but with randomly permuted section order. This enables testing whether organizational structure itself affects performance.

\subsubsection{Information Preservation Test (H-M1)}

To validate that format transformation preserves diagnostic information:

\begin{enumerate}
\item Generate (raw\_error, structured\_error) pairs for 500 samples across 8 error types
\item Use a judge to reconstruct original error details from structured format
\item Measure reconstruction accuracy for each field (line\_number, error\_type, error\_message, code\_context)
\item Success criterion: >95\% reconstruction accuracy across all fields
\end{enumerate}

\subsubsection{Format Comparison (H-M2)}

Paired comparison design to isolate structure effects:

\begin{enumerate}
\item 500 error instances formatted in both Structured and Scrambled conditions
\item Apply self-repair for each condition
\item Binary success/failure outcome per sample
\item Statistical test: McNemar's test for paired binary outcomes
\item Success criterion: Structured > Scrambled with p < 0.05
\end{enumerate}

\subsubsection{Fix Specificity Framework (H-M3)}

Four specificity levels based on scaffolding theory:

\begin{table}[htbp]
\centering
\resizebox{\columnwidth}{!}{%
\begin{tabular}{lll}
\toprule
Level & Name & Example \\
\midrule
0 & No hint & (diagnostic only) \\
1 & General strategy & ''convert types to match'' \\
2 & Specific pattern & ''use str(x) or int(y)'' \\
3 & Exact fix & ''str(count)'' \\
\bottomrule
\end{tabular}
}
\end{table}

The hypothesis predicts an inverted-U pattern with peak performance at intermediate levels (1-2).

\subsubsection{Scale Interaction Test (H-C1)}

2 (Format) $\times$ 3 (Model Scale) factorial design:

\begin{itemize}
\item Models: CodeLlama-7B-Instruct, CodeLlama-34B-Instruct, GPT-4
\item Analysis: Two-way ANOVA with interaction term
\item Success criterion: Interaction p < 0.05, $\eta^2$ > 0.01
\end{itemize}
